% SHARD-ID: CA-43-QUBIT-VACUUM-MOMENT-CONSTRAINTS
% SHARD-TITLE: Vacuum Moment Constraints for Qubit Residuals
% SHARD-SUMMARY: Defines finite restrictions of a translation-invariant global state by Pauli-word moments and positive moment matrices.
% SHARD-SUMMARY: Converts local Poincare and Virasoro residuals into linear moment equations of the form omega(X^* R Y)=0.
% SHARD-SUMMARY: Records the zero-expectation and normal-ordering data needed before an SDP can exclude Hamiltonians.
% SHARD-KEYWORDS: vacuum state, moments, GNS, Pauli words, SDP constraints, qubit chain

\section{Vacuum Moment Constraints for Qubit Residuals}
\label{sec:qubit-vacuum-moment-constraints}

\subsection{Status, Sources, and Dependencies}

\textbf{Status: proposal-level SDP interface backed by sourced state/GNS
architecture.}  The moment equations are finite algebra.  The existence of a
continuum vacuum or positive-energy representation is not claimed.

Dependencies: CA-15, CA-38--CA-42, CONVENTIONS.md (n).

% Source: literature/md/2010.11121/2010.11121.md:126--180 -- inductive-limit
%   algebra, projective-limit states, and GNS vector state.
% Source: literature/md/2010.11121/2010.11121.md:226--243 -- time-zero net,
%   vacuum invariance, positive energy, and Poincare-covariance warning.
% Source: references/text/CFTFromLatticeFermions.txt:376--388 -- vacuum /
%   ground-state subtraction in Koo-Saleur scaling.

\subsection{Finite Moment Data}

Let \(P_s\) be a finite Pauli word on \(\mathbb Z\), padded by identities
outside its support.  A candidate translation-invariant global state \(\omega\)
has moments
\[
  y_s=\omega(P_s),
  \qquad y_{\emptyset}=1,
  \qquad y_{\tau^n s}=y_s.
  \tag{43.1}
\]
For a finite word set \(\mathcal W\), positivity is the moment-matrix
constraint
\[
  M_{\mathcal W}(y)_{u,v}
  =
  \omega(P_u^*P_v)
  =
  \sum_s \mu_{u,v}^{s}y_s
  \succeq0,
  \tag{43.2}
\]
where the constants \(\mu_{u,v}^{s}\) are computed from the Pauli product
table of CA-38.  Hermiticity gives
\[
  M_{\mathcal W}(y)_{v,u}
  =
  \overline{M_{\mathcal W}(y)_{u,v}}.
\]
The SDP may be represented over real variables by splitting real and imaginary
parts.

\subsection{Zero Expectations}

The infinite sums \(H,P,K\) are not finite observables.  The finite data are
their local densities and residuals.  After choosing the scalar subtraction
\(\bar h=h-eI\), a vacuum-normalized test may impose
\[
  \omega(\bar h_j)=0,\qquad
  \omega(p_j)=0,
  \tag{43.3}
\]
and, for the residuals of CA-39,
\[
  \omega(A_j-Du_j)=0,\qquad
  \omega(B_j-u_j-v^2\bar h_j-Dw_j)=0.
  \tag{43.4}
\]
These are only the level-zero moment consequences.  They are usually too weak
to exclude much.

\subsection{GNS Relation Constraints}

The stronger finite restriction says that a residual vanishes on the vacuum
sector as a relation in the GNS representation.  For each chosen residual
density \(R(h)\) and probe words \(X,Y\),
\[
  \omega(X^*R(h)Y)=0.
  \tag{43.5}
\]
For fixed \(h\), fixed \(v^2,e,u,w\), and fixed residual coefficients, (43.5)
is linear in the moments \(y_s\).  The same rule applies to the finite
Witt/Virasoro residual densities \(r^{++}_{mn},r^{--}_{mn},r^{+-}_{mn}\) from
CA-42:
\[
  \omega(X^*r^{\pm\pm}_{mn}(h)Y)=0,
  \qquad
  \omega(X^*r^{+-}_{mn}(h)Y)=0.
  \tag{43.6}
\]
These equations are the promised finite data extracted from the demand that a
global state \(\omega\) make the generators and relations vanish in expectation
on a finite probe space.

\subsection{Highest-Weight Caveat}

If a chiral vacuum sector is named, one may also impose finite analogues of
\[
  \ell_n\Omega=0,\quad \bar\ell_n\Omega=0\qquad(n>0),
\]
for example
\[
  \omega(\ell_{-n}\ell_n)=0.
  \tag{43.7}
\]
But central charge, normal ordering, and finite-volume mode normalisation are
still convention-sensitive.  Until CONVENTIONS.md (d) is fixed for this block,
such constraints are optional diagnostics rather than canonical Virasoro
vacuum equations.

